
@@M0@@

Cette relation elle-même montre l'additivité de l'application @@M1@@, qui peut donc être prolongée en un homomorphisme de @@M2@@ dans @@M3@@~:

\textbf{Remarque 1.13.} On peut prendre pour base de @@M4@@ sur @@M5@@ n'importe quelle suite de puissances de @@M6@@~; @@M7@@, car @@M8@@ est inversible, et une telle suite est donc composée d'éléments linéairement indépendants sur @@M9@@; par ailleurs, les autres puissances de @@M10@@ peuvent s'exprimer explicitement comme combinaisons linéaires à coefficients dans @@M11@@ de celles-ci, d'après l'égalité (1.9.1), et l'inversibilité de @@M12@@.

Posons @@M13@@ (@@M14@@ s'interprète comme «classe d'une section hyperplane»). On peut alors prendre pour base de @@M15@@ sur @@M16@@ les éléments @@M17@@. Pour obtenir l'équation de dépendance de @@M18@@ on considère la suite exacte :
@@M19@@
et la suite exacte obtenue par passage aux duals :
@@M20@@
On remarque alors que @@M21@@, d'où l'on déduit :
@@M22@@
soit :
@@M23@@

On obtient un résultat analogue en introduisant @@M24@@.

Notons enfin que si @@M25@@, donc @@M26@@, on a @@M27@@, et l'égalité (1.13.1) se réduit à :
